% paper/sections/08f_pressure_summary.tex
%
% §8 章末まとめ：精度バランスと設計上のトレードオフ
% ※ 旧所在: 08c_ccd_poisson.tex（§8.2 内）→ §8 章末に移動（CHK-087: R5）
%   理由：GFM（§8.3）・BC（§8.4）を全て読んだ後に精度全体像を示す構成に変更
%% ═══════════════════════════════════════════════════════════════════════════════
\subsection{本章のまとめ：精度バランスと設計上のトレードオフ}
\label{sec:pressure_accuracy_summary}
%% ═══════════════════════════════════════════════════════════════════════════════

本章では，変密度 PPE の導出（§~\ref{sec:projection_derivation}），
CCD による $\Ord{h^6}$ 離散化（§~\ref{sec:ccd_poisson_matrix}），
欠陥補正法による効率的求解（§~\ref{sec:defect_correction_main}），
Extension PDE による界面越し場延長（§~\ref{sec:gfm}），
境界条件と数値安定性（§~\ref{sec:ppe_bc_stability}）を順に論じた．
以下に，各コンポーネントの精度を整理する．

\begin{table}[htbp]
\centering
\caption{数値手法コンポーネント別の精度まとめ}
\label{tab:accuracy_summary}
\renewcommand{\arraystretch}{1.3}
\footnotesize
\begin{tabular}{lp{0.28\linewidth}cc}
\toprule
コンポーネント & 手法 & 空間精度 & 時間精度 \\
\midrule
CLS 移流（界面追跡） & Dissipative CCD + TVD-RK3 & $\Ord{h^5}$ & $\Ord{\Delta t^3}$ \\
幾何量計算（曲率等） & CCD 法 & $\Ord{h^6}$ & — \\
PPE（離散化＋求解） & CCD 微分形式 + 欠陥補正法 + Extension PDE & $\Ord{h^6}$ & — \\
NS 粘性項 & CCD + Crank--Nicolson & $\Ord{h^6}$ & $\Ord{\Delta t^2}$ \\
NS 対流項（Predictor） & CCD（1次微分）+ AB2 外挿 & $\Ord{h^6}$ & $\Ord{\Delta t^2}$ \\
Projection スプリッティング & IPC 法（van Kan 1986） & — & $\Ord{\Delta t^2}$ \\
Extension PDE（界面越し場延長） & CSF + Extension PDE（Aslam 2004） & $\Ord{h^1}$（界面近傍†） & — \\
\midrule
\textbf{全体（空間律速：Extension PDE 界面延長†）} & & $\boldsymbol{\Ord{h^1}}$（界面近傍） & — \\
\textbf{全体（時間律速：AB2+IPC 整合）} & & — & $\boldsymbol{\Ord{\Delta t^2}}$ \\
\bottomrule
\end{tabular}

\smallskip
\noindent†\ Extension PDE（§~\ref{sec:extension_pde}）により界面直近の圧力場を滑らかに延長するが，
界面近傍（1–2格子幅以内）では Extension PDE の収束次数（一次精度，上風有限差分）が
全体の空間精度を律速する．平滑領域では全コンポーネント $\Ord{h^5}$ 以上が確保される．
\end{table}

\begin{tcolorbox}[warnbox, title={精度ミスマッチと設計上のトレードオフ}]
\begin{enumerate}
  \item \textbf{空間精度：}
        全数値コンポーネントが $\Ord{h^5}$--$\Ord{h^6}$ に揃っている（平滑領域）．
        Extension PDE（§~\ref{sec:extension_pde}）により界面直近の CCD Gibbs 振動を除去したが，
        Extension PDE 自体の一次精度（上風 FD）が界面近傍の律速段階となる．

  \item \textbf{時間精度：}
        AB2 + IPC により全体 $\Ord{\Delta t^2}$（§~\ref{sec:ipc_derivation}）．
        対流（AB2），粘性（Crank--Nicolson），Projection（IPC）の3要素が統一されている．
\end{enumerate}
\end{tcolorbox}
